\section{Discussion and Conclusion}
\label{sec:conclusion}

\paragraph{Discussion.}
The picture across SWE-Bench Verified
(Table~\ref{tab:verified}) and SWE-PolyBench
(Table~\ref{tab:swepoly-java-python},
Figure~\ref{fig:scorer-func-recall}) is consistent. First, a
bounded-context scaffolding closes most of the gap left by the
first-stage retriever: every LLM-driven configuration we tried
exceeds the $\approx0.35$ Recall@5 ceiling of the zero-shot
ranker. Second, graph exploration (Step~B) is the largest single
lever, lifting recall and accuracy on every model and language
while keeping per-call context short. Third, a small fine-tuned
scorer is enough to drive the full scaffolding: a
Qwen2.5-Coder-3B scorer matches Qwen3-Coder-30B-A3B under the
same scaffolding on Verified at roughly an order of magnitude
lower token cost ($35.6$k vs $\geq 198$k tokens per issue), and
pairs with a 30B Step-C agent to set the best Python numbers on
SWE-PolyBench.
Token consumption tracks the call count rather than the model
size, so larger Step-C agents do not pay a context tax for using
our scorer.

\paragraph{Key takeaways.}
(i) Localization scales better through scaffolding structure
than through model size: a 3B scorer plus a 30B agent under our
loop beats a 30B agent alone. (ii) Graph-aware exploration is
cheap because the scorer sees short neighborhood windows rather
than the full pool; the cost of an extra step is small compared
to the recall it recovers. (iii) Supervised training is
sufficient for the scorer because the per-round shaping signal
telescopes to end-to-end Recall@$K$, so no policy-gradient
machinery is needed for the inner loop.

\paragraph{Conclusion.}
We cast code localization as sequential set-coverage on a code
graph and built a bounded-cost, multi-round, multi-agent
scaffolding around that view. The round structure exposes a
dense per-round target whose optimum coincides with end-to-end
Recall@$K$, which lets a small open scorer be trained by
ordinary supervised learning rather than sparse-reward RL. The
resulting Qwen2.5-Coder-3B scorer makes the full scaffolding
competitive with GPT-5.3 Codex on SWE-Bench Verified and on
SWE-PolyBench Java and Python, at a fraction of the token cost.
Stop-timing and cross-round coordination remain natural
candidates for a light-weight RL stage on top of the supervised
scorer; we leave that to future work.
